\chapter{Sprint 64 --- Cụm G: Issue List lứa 140--151, G1 $\to$ G6 (2026-09-26 $\to$ 09-30)}

\emph{Chapter này ghi lại cụm G: lứa issue BA mở trên tab \texttt{5. Issue List} sau khi khối A (chapter 77) khép.
Mười hai issue (11 có code, 146 tạm dừng), bảy lượt code, một lượt review trên UAT. Mỗi lượt lên UAT ngay trong ngày, đo lại ở chế độ chỉ-đọc rồi
mới chuyển sheet sang \emph{Ready for Retest}. Dev không tự đóng \emph{Done}. Nhật ký ở \texttt{docs/f-progress.md}; prompt
và bài học từng lượt ở \texttt{docs/next-session-clusters-G.md}; nghiệm thu ở \texttt{docs/g1-} \ldots\
\texttt{g6-acceptance-matrix.md}; review UAT ở \texttt{docs/g-review-uat.md}.}

\section{Vì sao có cụm G}

Ngày 26/09, \texttt{issue\_queue.py --dev} trả 7 issue \emph{Ready for Dev} (STT 140--146). Reconcile bằng
\texttt{git log --all -S"<ID>"} và \texttt{grep} trong \texttt{custom/} cho cả 7: chỉ có dòng nhắc trong docs, 0 dòng code.
Vậy cả 7 đều phải làm thật, không có lượt ``chỉ ghi ledger''. Chủ dự án chốt ba điều:
\begin{enumerate}[nosep]
\item Chuẩn hoá component trước, trang sau: G1 component mobile $\to$ G2 shell $\to$ G3 Home PC $\to$ G4 routing.
\item 144 làm cuối G1, đo sau khi 143 đã đổi token (BA ghi 144 phụ thuộc 143).
\item Mỗi cụm một phiên, một lần \texttt{-u}, một lần deploy; cụm lớn (G3) tách a/b.
\end{enumerate}
Ngày 29/09, BA chạy exploratory trên UAT bằng đơn test S00075 và mở thêm 147--151. Cùng lúc, dòng 146 biến khỏi sheet.
Hai issue Home (150, 151) làm ngay thành lượt H150/151. Ba issue Đặt hàng/Lịch sử đánh số tiếp G5, G6.

\begin{center}
\footnotesize
\begin{tabularx}{\textwidth}{>{\raggedright\arraybackslash}p{1.6cm}>{\raggedright\arraybackslash}p{4.3cm}>{\raggedright\arraybackslash}p{1.5cm}>{\raggedright\arraybackslash}X}
\toprule
Lượt & Issue (STT) & Commit & Module (version sau lượt) \\
\midrule
G1 & \texttt{UI-MOB-HEADER-DENSITY-001} (143), \texttt{UI-MOB-BOTTOMNAV-DENSITY-001} (145), \texttt{WJ-ORD-MOB-SPACING-001} (144) &
  \texttt{fcce811} & \texttt{portal\_layout} 19.0.59.0.0 · \texttt{portal\_exam} 19.0.6.1.0 · \texttt{portal\_sale} 19.0.4.26.0 \\
G2 & \texttt{UI-MOB-STORE-SWITCHER-001} (141), \texttt{UI-PC-TOPBAR-REG-001} (140) &
  \texttt{6745671} & \texttt{portal\_base} 19.0.7.29.0 · \texttt{portal\_layout} 19.0.59.1.0 \\
G3a & \texttt{UI-PC-HOME-REDESIGN-001} (142) --- khung &
  \texttt{160d13e} & \texttt{portal\_base} 19.0.7.30.0 · \texttt{portal\_debt} 19.0.4.16.0 \\
G3b & \texttt{UI-PC-HOME-REDESIGN-001} (142) --- 7 block &
  \texttt{d8f89bf} & \texttt{portal\_base} 19.0.7.31.0 \\
H150/151 & \texttt{WJ-HOME-009} (150), \texttt{WJ-HOME-010} (151) &
  \texttt{0cd1f0a} & \texttt{portal\_base} 19.0.7.32.0 · \texttt{purchase\_history} 19.0.3.20.0 \\
G4 & \texttt{WJ-PORTAL-ROUTING-001} (146) & --- & Tạm dừng: BA xoá dòng khỏi sheet, 0 dòng code \\
G5 & \texttt{WJ-ORD-028} (148), \texttt{WJ-ORD-027} (147) &
  \texttt{887da0a} & \texttt{portal\_layout} 19.0.59.2.0 · \texttt{portal\_base} 19.0.7.33.0 · \texttt{portal\_sale} 19.0.4.27.0 \\
G6 & \texttt{WJ-ORD-029} (149) &
  \texttt{9efa67f} & \texttt{purchase\_history} 19.0.3.21.0 \\
\bottomrule
\end{tabularx}
\end{center}

Cuối sprint, cả 11 issue có code đều ở \emph{Ready for Retest} trên sheet, đã đọc lại bằng CSV.
\texttt{issue\_queue.py --dev} ngày 01/10 trả 0.

\section{G1 --- Mật độ mobile: PageHeader, SectionHeader, BottomNav, nhịp Đặt hàng}

BA đo trên UAT thấy đầu trang mobile chiếm nhiều chỗ. Sửa ở token dùng chung, nên mọi trang portal mobile đều được:
\begin{itemize}[nosep]
\item \textbf{143}: PageHeader mobile 52 $\to$ \textbf{44} cho cả 3 kiểu (tiêu đề / có nút lùi / có nút tạo). Chủ dự án chốt
  ba kiểu cùng chiều cao. SectionHeader mobile còn 18/24px, chỉ đổi trong \texttt{@media} mobile.
\item \textbf{145}: thanh điều hướng dưới 83 $\to$ \textbf{72 + vùng an toàn}, mỗi mục 50, nhãn 12, badge neo góc icon.
  Sửa kèm công thức \texttt{--wujia-mnav-total} sai từ trước: sheet ``Thêm'' và nút nổi của màn Thi giờ bám đúng mép
  trên thanh, kể cả khi máy có vùng an toàn.
\item \textbf{144} (làm cuối): trang Đặt hàng mobile, ô tìm $\to$ chip 23 $\to$ \textbf{12}px, chip $\to$ ``Danh sách sản phẩm''
  14 $\to$ \textbf{8}px. Khoảng 23 cũ đến từ \texttt{margin} 15 của \texttt{form} toàn cục, nên chỉ sửa trong
  \texttt{.wujia-morder}.
\end{itemize}
\textbf{Số đo}: 28 route $\times$ 360/390/430 $\times$ có/không vùng an toàn 34: PageHeader 44 ở mọi kiểu, 0 tràn ngang,
cuộn tới cuối vẫn thấy nút cuối (khe $\geq$13). PC không đổi (vân tay 76/81, 5 chỗ khác do bộ đếm lượt xem). 16 test tĩnh,
mutation 11/11, suite 930/0/0. Bộ đo mới \texttt{scripts/qa/wj\_density.py} (Playwright, cờ \texttt{--safe-area},
\texttt{--readonly}) về sau được dùng lại cho mọi lượt có đụng mobile.

\section{G2 --- Dải cửa hàng mobile (141) và top bar PC (140)}

\begin{itemize}[nosep]
\item \textbf{141}: dải ``cửa hàng đang thao tác'' hiện trên mọi trang mobile, kể cả Home. Có chevron \emph{chỉ} khi user
  có hơn một cửa hàng; user một cửa hàng thấy dải tĩnh. Bấm vào bất kỳ chỗ nào trên dải đều mở hộp chọn cửa hàng.
\item \textbf{Nhãn vai trò một nguồn tiếng Việt}: Chủ tiệm / Quản lý / Nhân viên. \texttt{ROLE\_LABELS} dời về model
  \texttt{wujia.franchise.member} (\texttt{\_portal\_role\_label()}) và dùng ở 5 chỗ: dải mobile, khối Cửa hàng PC,
  menu tài khoản, chữ dưới avatar, hero Home mobile.
\item \textbf{140}: icon giỏ/chuông trên top bar PC lệch tâm khoảng 9px từ cụm B (04/08). Nguyên nhân là
  \texttt{.nav-link\{display:block\}} của \texttt{web.assets\_frontend} nạp sau. Nay icon nằm giữa vòng tròn 40, badge
  không che icon với 0/4/12, chip vai trò nằm trong khối Cửa hàng.
\end{itemize}
\textbf{Số đo}: lệch tâm icon ($-$9.4, $-$9.5) $\to$ (0, 0); phần badge chồng lên icon 59.4 $\to$ 0 px². Mobile 27 route
$\times$ 3 khổ $\times$ 2 user không tràn. 11 test, mutation 9/9, suite 941/0/0.

\section{G3a/G3b --- Home PC theo mockup V4 (142)}

Home PC cũ có hero và ``Thao tác nhanh''; mobile thì đã theo bố cục mới. 142 đưa PC về cùng nội dung với mobile theo
mockup V4. Chia hai lượt, deploy gộp một lần để BA không thấy Home nửa cũ nửa mới.
\begin{itemize}[nosep]
\item \textbf{G3a --- khung}: hàng đầu chia 50/50 (card Cửa hàng hiện tại | card Khung giờ đặt hàng có thanh tiến độ),
  4 KPI Đơn hàng · Thông báo · Đổi trả · Công nợ. Các KPI dùng cùng biến và link với mobile. Bỏ KPI ``Đơn chờ xử lý'' và
  bảng ``Sản phẩm mua nhiều nhất'' theo danh sách BA. Ô Công nợ có số thật qua khe \texttt{home\_debt\_kpi} của
  \texttt{portal\_debt}; nếu gỡ module thì ô hiện ``---''. Query \texttt{/portal} \textbf{$-$7} (46 $\to$ 39).
\item \textbf{G3b --- 7 block record}: Đơn hàng · Đổi trả · Giao hàng · Thông báo · Kiến thức · Hỗ trợ · Thông tin cửa
  hàng. Dòng chép từ block mobile (cùng biến, cùng tiền, cùng badge, cùng link). Lưới CSS \textbf{2 cột ở 992--1399, 3 cột
  từ 1400}; ngưỡng đổi từ 1200 sau khi đo ra sidebar chiếm chỗ từ 1200. Màu loại thông báo gom về một
  \texttt{noti\_badge\_map} cho cả hai kênh. Controller không sửa; query $\Delta$0 so với G3a.
\end{itemize}
\textbf{Số đo}: 7 khổ $\times$ 3 user, kể cả chuỗi dài VI/EN/ZH: 0 tràn, 0 chữ bị cắt, 0 badge đè; các card trong cùng hàng
cao bằng nhau ($\Delta h$ 0). Mobile $\Delta$0. 34 test, mutation 7/7 + 7/7, suite 975/0/0. Nghiệm thu 12/12.

\section{H150/151 --- Home giống Lịch sử}

Cùng đơn S00075, BA thấy Home PC in \emph{15:57} còn Lịch sử in \emph{22:57} (150), và Home in ``Nháp'' còn Lịch sử in
``Chờ xác nhận'' (151).
\begin{itemize}[nosep]
\item \textbf{150} đã hết từ G3b: block PC cũ in \texttt{date\_order} giờ UTC; G3b thay bằng \texttt{wj\_dt} (múi giờ người
  dùng). Lượt này chỉ thêm test chặn hồi quy.
\item \textbf{151}: Home có bảng nhãn riêng \texttt{MOBILE\_ORDER\_BADGES} (draft $\to$ ``Nháp''). Luật trạng thái đơn của Lịch
  sử dời xuống \texttt{wujia\_portal\_base/controllers/utils.py} (\texttt{portal\_order\_status}, \texttt{portal\_order\_badge});
  Home PC, Home mobile và Lịch sử cùng gọi. Đã xoá bảng nhãn riêng.
\end{itemize}
\textbf{Số đo}: 4 đơn trên UAT trùng chữ, màu và giờ ở Home PC, Home mobile và chi tiết Lịch sử. Query $\Delta$0.
9 test, mutation 6/6.

\section{Review cụm G trên UAT}

Lượt review chỉ đọc sau H150/151, chạy lại toàn bộ bộ đo ở chế độ chặn ghi, \textbf{0 request bị chặn}. Version UAT
khớp HEAD ở cả 6 module. 25 dòng AC: \textbf{20 đạt · 1 không đạt · 4 LIMIT}.
\begin{itemize}[nosep]
\item \textbf{Không đạt}: top bar PC vỡ ở \emph{đúng} 992px. Nút hamburger nhận \texttt{margin-right} 451.7px nên khối Cửa
  hàng rớt xuống hàng dưới và bị header cắt. Từ 993 trở lên bình thường. Lỗi có từ trước G2.
\item \textbf{LIMIT}: UAT không có user nhiều cửa hàng (141 đã đạt trên DB copy); nav \texttt{em.hcm} không có badge lúc đo;
  bảng giá \texttt{Default} trên UAT là USD nên đơn portal in \texttt{\$}. Đây là cấu hình, không phải code.
\end{itemize}

\section{G4 --- Điều hướng \texttt{/} (146): tạm dừng}

146 xin: vào \texttt{/} khi chưa đăng nhập thì sang \texttt{/portal/login}; user nội bộ sang \texttt{/web}; cửa hàng sang
\texttt{/portal}. Đã soi xong và chủ dự án đã chốt cách làm. Override \texttt{Home.index} không có tác dụng, vì lá controller
của \texttt{website} đứng trước trong MRO. Cách chốt là kế thừa \texttt{ir.http.\_pre\_dispatch} trong
\texttt{wujia\_portal\_base}, không depend \texttt{website}. Trước khi viết dòng code nào thì BA xoá dòng 146 khỏi sheet.
Cách làm giữ nguyên ở khối G4 của \texttt{next-session-clusters-G.md}, chờ BA mở lại.

\section{G5 --- Đặt hàng: xác nhận trước khi gửi (148) và lọc gửi lại được (147)}

\begin{itemize}[nosep]
\item \textbf{148} (High): bấm ``Gửi đơn đặt hàng'' giờ mở hộp tóm tắt (số mặt hàng, tổng số lượng, tổng tiền, cửa hàng,
  ghi chú) với hai nút Hủy | Xác nhận gửi đơn, trên cả PC và mobile; phím Esc tương đương Hủy. Bấm Hủy thì giỏ giữ nguyên,
  không tạo gì. Chống tạo trùng mà không đổi schema: khoá nút + overlay, server giữ khoá \texttt{NOWAIT}. Nếu lần gửi thứ hai
  gặp giỏ rỗng mà vừa có đơn draft/sent của đúng cửa hàng và user trong 2 phút thì chuyển sang màn kết quả của đơn đó.
\item \textbf{147}: tái hiện trước rồi mới sửa. Gốc lỗi không nằm ở trang Đặt hàng. Khoá chống bấm đúp dùng chung
  (CMP-BTN-001, E6a) chỉ nhả khi tải lại trang, nên mọi form lọc bằng AJAX chỉ gửi được một lần. Lỗi này dính
  \textbf{11 màn danh sách}, BA chỉ thấy ở Đặt hàng. Sửa bằng sự kiện \texttt{wj:form:release} trong
  \texttt{wj\_ajax\_list.js}.
\item \textbf{Sửa kèm}: nút ``Gửi đơn đặt hàng'' trên mobile không gửi được từ E6a. Listener overlay tự kiểm lại cờ mà khoá
  chung đã cắm trước nên tự chặn lần gửi đầu.
\end{itemize}
\textbf{Số đo}: 13 test, mutation 12/12, suite 1006 (chỉ còn error fixture \texttt{wujia\_sale} có sẵn), query $\Delta$0 trên
27 route. UAT chỉ-đọc: hộp xác nhận 8/8 (chỉ mở hộp rồi Hủy, 0 request gửi đơn), lọc gửi lại 21/21. Nghiệm thu 148 6/6,
147 4/4.

\section{G6 --- Lịch sử: ``Ngày xác nhận'' chỉ khi đơn đã xác nhận (149)}

Đơn S00075 đang ``Chờ xác nhận'' nhưng trang chi tiết vẫn in ``Ngày xác nhận 29/09 22:57''. Nguyên nhân: Odoo dùng
\texttt{date\_order} làm giờ tạo báo giá và chỉ ghi đè lúc xác nhận. Sửa trong \texttt{\_history\_row\_vals} /
\texttt{\_history\_detail\_vals} bằng một key mới \texttt{confirm\_date}, chỉ có giá trị khi \texttt{state == 'sale'}
(Odoo 19 không có \texttt{done}). Bảng PC in ``---''; chi tiết ẩn hẳn dòng; nhãn ``Ngày tạo'' đổi thành
``Ngày đặt hàng'' (vẫn là \texttt{create\_date}). Lọc theo ngày vốn đã dùng \texttt{create\_date} nên không đổi.
\textbf{Số đo}: 7 test, mutation 7/7, 393/0/0, query $\Delta$0; UAT chỉ-đọc 12/12; nghiệm thu 5/5.

\section{Gì đã làm (nghiệp vụ)}

\begin{itemize}[nosep]
\item \textbf{Mobile gọn hơn}: đầu trang, tiêu đề mục và thanh điều hướng dưới đều thấp hơn, nên màn hình có thêm chỗ
  cho nội dung. Trang Đặt hàng bớt thưa ở phần đầu.
\item \textbf{Luôn biết đang thao tác cửa hàng nào}: dải cửa hàng có ở mọi trang mobile; vai trò ghi tiếng Việt thống nhất
  ở mọi chỗ.
\item \textbf{Home PC mới}: cùng nội dung với mobile. Nhìn một lần thấy cửa hàng, khung giờ, 4 chỉ số và 7 khối việc
  gần đây. Giờ và trạng thái đơn trên Home khớp hẳn với Lịch sử.
\item \textbf{Đặt hàng an toàn hơn}: có bước xác nhận trước khi gửi, không tạo nhầm hay tạo trùng đơn. Tìm/lọc đổi
  điều kiện bao nhiêu lần cũng được, trên cả 11 màn danh sách. Nút gửi đơn trên mobile hoạt động lại.
\item \textbf{Lịch sử không còn mâu thuẫn}: đơn chưa xác nhận thì không có ``Ngày xác nhận''.
\end{itemize}

\section{Lý do / trade-off}

\begin{itemize}[nosep]
\item \textbf{Sửa ở token/component, không sửa từng trang} (G1, G2, 147): một chỗ sửa ăn cả portal. Đổi lại, lượt nào cũng
  phải đo đủ 26--28 route $\times$ 3 khổ để chắc không trang nào vỡ; cái giá này đã có bộ đo dùng lại được.
\item \textbf{148 không thêm field/bảng}: chống trùng bằng khoá nút + khoá \texttt{NOWAIT} + chuyển về đơn vừa tạo.
  Không cần migration và không đụng quyền. Giới hạn: cửa sổ ``vừa tạo'' cố định 2 phút.
\item \textbf{Key mới \texttt{confirm\_date} thay vì đổi nghĩa \texttt{date\_order}} (G6): caller khác (Home) vẫn dùng
  key cũ không vỡ.
\item \textbf{Deploy G3a + G3b một lần}: BA không bao giờ thấy Home PC nửa cũ nửa mới, đổi lại G3a nằm trên \texttt{main}
  vài giờ mà chưa lên UAT.
\item \textbf{Lưới CSS thay \texttt{col-*} Bootstrap} (G3b): BS4 của Vuexy và BS5 của Odoo xung đột thứ tự nạp; CSS grid
  không phụ thuộc thứ tự.
\end{itemize}

\subsection*{Bài học}
\begin{itemize}[nosep]
\item \textbf{Runner deploy có thể bỏ lỡ một lượt push} (2 lần: G3a, G5). Luôn đọc version module trên UAT qua XML-RPC,
  đừng tin ``đã push''.
\item \textbf{Bridge sheet trả lỗi không phải JSON dù đã ghi}. Phải đọc lại CSV trước khi ghi lại; chạy lại
  \texttt{qa\_sync --apply} sẽ ghi History sai. Kèm theo: \texttt{qa\_sync} đã sửa sang dò cột theo tên header, vì BA
  chèn cột R và tool cũ ghi đè R bằng số cột cứng.
\item \textbf{Lỗi chỉ ở một khổ đúng mốc} (992): quét $\pm$1px quanh mỗi breakpoint, rồi so computed style hai khổ để ra
  ngay thuộc tính gây lỗi.
\item \textbf{Mọi nhãn trạng thái đơn đi qua \texttt{portal\_order\_status}}. Bảng nhãn ``riêng cho mobile'' là nguồn lệch.
\item \textbf{Khoá dùng chung phải có đường nhả} cho trang gửi bằng JS mà ở lại trang (\texttt{wj:form:release}). Đừng để
  lớp sau kiểm lại cờ của lớp trước.
\item \textbf{Issue có thể biến khỏi sheet} mà \texttt{issue\_queue --dev} vẫn báo 0. Đầu phiên phải so STT cuối sheet với
  bảng cụm đang làm.
\item \textbf{Đo ``gửi lại'' phải so URL}, không chỉ đếm request (route fragment \texttt{<path>/results}).
\end{itemize}

\subsection*{Nợ còn lại (chỉ ghi)}
\begin{itemize}[nosep]
\item Top bar PC vỡ ở đúng 992px (\texttt{wujia\_portal\_layout}, có từ trước G2); topbar 992--1199 chật, BA chưa tách issue.
\item Bảng giá \texttt{Default} trên UAT là USD (công ty VND): chủ dự án đổi cấu hình.
\item G4 (146) chờ BA xác nhận xoá hay chuyển chỗ.
\item 2 error fixture ở \texttt{wujia\_sale} (\texttt{test\_06\_filters\_return\_right\_orders}, setUpClass
  \texttt{TestWujiaSupplyDemandReport}).
\item Cột R của 6 dòng cũ (126, 129, 131, 132, 136, 139) bị \texttt{qa\_sync} cũ ghi đè, cần lấy lại qua Version history.
\item Nhãn vai trò chưa dịch EN/ZH; vùng an toàn chưa đo trên iPhone thật.
\item Dọn server đo 8032/8033/8055 + DB \texttt{wujia\_g2s}, \texttt{wujia\_g3s}, \texttt{wujia\_g5s} + \texttt{git worktree prune}.
\item Còn từ chapter 77: \texttt{auto\_install} theo ADR-027 · 7 câu hỏi BA ADR-027 · bàn giao anh Thái.
\end{itemize}

\textbf{Sau cụm G}: hàng đợi Dev trống (01/10). BA retest 140--145, 147--151. Phiên kế lấy issue mới khi BA mở,
hoặc làm nợ top bar 992.
